\section{Evaluation Metrics}

\begin{table}[H]
\centering
\begin{tabular}{llll}
\toprule
\textbf{Metric} & \textbf{Meaning} & \textbf{RAG v1.0} & \textbf{Agent v2.0}\\
\midrule
Recall@5 (PATH 1)   & Top-5 correct category/color   & $\geq$85\% & $\geq$88\%\\
MRR                 & Position of first correct result & $\geq$0.75 & $\geq$0.80\\
Hit Rate@5          & $\geq$1 correct result in top-5 & $\geq$90\% & $\geq$93\%\\
Task Completion     & No crash, returns results       & N/A         & $\geq$95\%\\
Faithfulness        & No hallucinated information     & N/A         & $\geq$0.87\\
Latency P95         & Total ReAct loop time           & $<$15s      & $<$15s\\
\bottomrule
\end{tabular}
\caption{System evaluation metrics comparing RAG v1.0 and Agent v2.0}
\end{table}

\section{Data Collection and Evidence Integrity}
\label{sec:data_evidence_integrity}

To make the final conclusions auditable, this thesis reports not only model metrics but also the provenance of collected data, aggregation boundaries, and data-quality checks. The evidence pipeline is based on operational PostgreSQL logs generated by the deployed system.

\subsection{Canonical Data Sources}
\label{subsec:canonical_data_sources}

\begin{table}[H]
\centering
\begin{tabular}{p{3.2cm}p{4.1cm}p{6.6cm}}
\toprule
\textbf{Evidence Signal} & \textbf{Source Table / View} & \textbf{Purpose in Evaluation}\\
\midrule
LLM token usage & \texttt{llm\_token\_usage}, \texttt{session\_token\_summary}, \texttt{mode\_cost\_summary} & Measure input/output token consumption per session and estimate operational cost by orchestration mode.\\
Behavior funnel & \texttt{product\_impressions}, \texttt{product\_clicks}, \texttt{selected\_items}, \texttt{product\_intents}, \texttt{user\_orders}, \texttt{session\_path\_funnel\_summary} & Quantify user progression from impression to selection (cart add), purchase intent, and order conversion.\\
Retrieval quality & Evaluation query set + top-\textit{k} ranked outputs & Compute Recall@5, MRR, Hit Rate@5, and faithfulness-based quality indicators.\\
\bottomrule
\end{tabular}
\caption{Canonical evidence sources used for thesis conclusions}
\label{tab:canonical_evidence_sources}
\end{table}

\subsection{Metric Definitions and Cohort Boundaries}
\label{subsec:metric_definitions_boundaries}

\textbf{Token metrics.}
For a cohort of sessions $\mathcal{S}$ with valid synthesis records (\texttt{call\_name='synthesis'}), token usage is aggregated as:
\begin{equation}
T_{in}=\sum_{s \in \mathcal{S}} \text{input\_tokens}_s,\quad
T_{out}=\sum_{s \in \mathcal{S}} \text{output\_tokens}_s,\quad
T_{total}=T_{in}+T_{out}
\end{equation}
When orchestration is enabled, orchestrator token fields are reported separately through \texttt{mode\_cost\_summary} to prevent hidden cost inflation.

\textbf{Add-to-cart and behavior metrics.}
In this system, a cart-add event corresponds to a row in \texttt{selected\_items}. Let $I$ = impressions, $C$ = clicks, $A$ = cart-add selections, $W$ = ``will\_buy'' intent, $O$ = orders:
\begin{equation}
\text{CartAddRate}=\frac{A}{I},\quad
\text{WillBuyRate}=\frac{W}{A},\quad
\text{ConversionRate}=\frac{O}{|\mathcal{S}|}
\end{equation}
These definitions align with analytics endpoints in \texttt{api/analytics.py} and the funnel view in \texttt{agent/memory.py}.

\textbf{Accuracy metrics.}
Retrieval quality is reported using Recall@5, MRR, and Hit Rate@5 (Table~\ref{tab:accuracy_reporting_template}), computed on an evaluation cohort of $N_{query}$ labeled queries. Every reported score must include $N_{query}$ to avoid over-interpreting small samples.

\begin{table}[H]
\centering
\begin{tabular}{p{2.8cm}p{8.4cm}p{2.4cm}}
\toprule
\textbf{Metric} & \textbf{Definition} & \textbf{Report with}\\
\midrule
Recall@5 & Proportion of queries where at least one relevant item appears in top-5 results. & $N_{query}$\\
MRR & Mean reciprocal rank of the first relevant retrieved item. & $N_{query}$\\
Hit Rate@5 & Fraction of queries with $\geq$1 relevant result in top-5. & $N_{query}$\\
\bottomrule
\end{tabular}
\caption{Accuracy reporting template with mandatory sample-size context}
\label{tab:accuracy_reporting_template}
\end{table}

\subsection{Inclusion, Exclusion, and Time-Window Rules}
\label{subsec:inclusion_exclusion_rules}

To ensure reproducibility, all aggregates in this chapter SHALL follow the same boundary policy:
\begin{enumerate}
    \item Include only sessions inside the declared collection window $[t_{start}, t_{end}]$.
    \item Exclude sessions with missing key fields for the target metric (e.g., missing token fields for token analysis).
    \item Exclude duplicate interaction rows when the same session-event key appears multiple times due to retries.
    \item Report both the original cohort size and the post-filter size to make exclusion impact explicit.
\end{enumerate}

\subsection{Data Quality and Validity Checks}
\label{subsec:data_quality_validity}

\begin{table}[H]
\centering
\begin{tabular}{p{4.0cm}p{4.7cm}p{5.0cm}}
\toprule
\textbf{Check} & \textbf{Validation Rule} & \textbf{Interpretation if Failed}\\
\midrule
Missingness & Required fields are non-null (session\_id, event timestamp, metric-specific columns). & Report as undercount risk; do not claim strong support from affected metric.\\
Consistency & Event semantics are stable (e.g., cart-add is always \texttt{selected\_items}). & Re-map metric definitions before comparing across sections.\\
Duplicate events & No duplicated session-event records after deduplication policy. & Treat inflated rates as invalid until corrected.\\
Sample balance & Cohort not dominated by one subgroup/mode. & Downgrade claim strength to ``partially supported'' or ``inconclusive.''\\
\bottomrule
\end{tabular}
\caption{Validity checks for deciding whether collected data is trustworthy}
\label{tab:data_validity_checks}
\end{table}

This validity layer is the decision gate for ``data right or wrong'': metrics are considered trustworthy only when they pass the checks in Table~\ref{tab:data_validity_checks}; otherwise, conclusions must be weakened accordingly.

\section{Design Strengths}
\begin{enumerate}
    \item \textbf{Intent-first}: The agent classifies intent before acting --- prevents incorrect searches.
    \item \textbf{Clarification gate}: For ambiguous queries, asks instead of guessing --- improves quality.
    \item \textbf{Graceful degradation}: Politely refuses when no product is found, with alternative suggestions.
    \item \textbf{Hybrid retrieval}: BM25 + FashionSigLIP + RRF leverages the strengths of both approaches.
    \item \textbf{PostgreSQL}: Persistent, scalable, supports concurrent access better than CSV.
    \item \textbf{Semantic token (caption)}: LLM-generated descriptions enable deeper embedding understanding beyond short labels.
    \item \textbf{Query Expansion}: Gemini-powered synonym queries increase recall for short queries.
    \item \textbf{Cross-encoder Reranker}: BGE v2-m3 reranks results for higher precision.
\end{enumerate}

\section{Risks and Mitigation}

\begin{table}[H]
\centering
\begin{tabular}{lll}
\toprule
\textbf{Risk} & \textbf{Severity} & \textbf{Mitigation}\\
\midrule
High latency (multiple Gemini calls) & High   & Parallel tool calls, cached expansion\\
LLM hallucination                    & High   & RAGAS faithfulness $\geq$0.87, fallback templates\\
BM25 rebuild on restart              & Low    & Serialize BM25 index to disk\\
PostgreSQL cold start                & Low    & Connection pooling\\
Gemini API rate limit                & Medium & Batch processing, exponential backoff\\
\bottomrule
\end{tabular}
\caption{Risk matrix and corresponding mitigation strategies}
\end{table}

\section{Conclusion}
This project successfully built an intelligent multimodal fashion search system following a clear two-phase roadmap. The RAG v1.0 phase established a hybrid search foundation with FashionCLIP 2.0 and BM25, while Agent v2.0 PATH 1 upgraded to a proactive architecture featuring a ReAct loop, intent classification, a clarification gate, and persistent memory via PostgreSQL.

Version v2.0 represents significant improvements over v1.0:
\begin{itemize}
    \item \textbf{Embedding}: Upgraded from FashionCLIP (512-d) to Marqo-FashionSigLIP (768-d).
    \item \textbf{Query Expansion}: Added Gemini-powered synonym expansion for short queries.
    \item \textbf{Reranker}: Integrated BGE Reranker v2-m3 cross-encoder for improved precision.
    \item \textbf{Deployment}: Docker Compose stack with 4 services (PostgreSQL + Qdrant + FastAPI + Cloudflare Tunnel).
\end{itemize}

\subsection{Evidence-Based Claim Mapping for Final Conclusion}
\label{subsec:evidence_claim_mapping}

Final claims are mapped to evidence strength so that conclusions remain proportional to verified data:

\begin{table}[H]
\centering
\begin{tabular}{p{4.0cm}p{5.8cm}p{3.0cm}}
\toprule
\textbf{Claim} & \textbf{Primary Evidence} & \textbf{Status}\\
\midrule
Agent v2.0 improves retrieval quality over v1.0. & Recall@5, MRR, Hit Rate@5 (Section~\ref{sec:data_evidence_integrity} and Section~\ref{subsec:metric_definitions_boundaries}). & Supported\\
Agent behavior is robust in multi-turn interaction. & Task completion, faithfulness, and clarification behavior logs. & Supported\\
Operational LLM cost is explainable and monitorable. & \texttt{llm\_token\_usage}, \texttt{session\_token\_summary}, mode cost aggregation. & Partially supported\\
User purchase tendency is improved by recommendations. & CartAddRate / WillBuyRate / ConversionRate from behavior funnel. & Inconclusive (depends on cohort size and balance)\\
\bottomrule
\end{tabular}
\caption{Claim-to-evidence mapping used in the final thesis conclusion}
\label{tab:claim_evidence_mapping}
\end{table}

Therefore, the final conclusion follows three rules: (1) only claims marked \textit{Supported} are stated as confirmed findings, (2) \textit{Partially supported} claims are reported with operational caveats, and (3) \textit{Inconclusive} claims are retained as hypotheses for future data collection.

\textbf{Current evidence limitations.}
\begin{itemize}
    \item Token-to-cost estimates depend on fixed public list prices and may vary over time.
    \item Behavior metrics are sensitive to sample imbalance across orchestration modes and user groups.
    \item Conversion-style claims require larger longitudinal cohorts than early-stage thesis collection windows.
\end{itemize}

\textbf{Future development directions:}
\begin{itemize}
    \item \textbf{Phase 3}: Scale dataset to 15K--100K products, fine-tune the reranker, Redis session cache.
    \item \textbf{Phase 4}: Virtual Try-on (IDM-VTON), fine-tune FashionSigLIP on proprietary data, A/B testing.
    \item \textbf{PATH 2}: Image-to-Image search with Composed Search.
\end{itemize}
